\chapter{Blocks 7--8: Statistics for Genomics}\label{chap:stats}

\textbf{Primary:} Distributions, multiple testing, linear models, PCA\quad|\quad\textbf{Interleave:} Python + R implementation, genomic data

\begin{brainbox}[Why this unit matters most]
Many weak genomics papers have a statistics problem, not a coding problem.

You do not need to derive proofs. You do need to know: which distribution fits my data? What test is appropriate? What happens when my assumptions are wrong?

This block teaches by simulation: generate data you already know the truth about, break assumptions on purpose, and see what the methods do.
\end{brainbox}


\section{Block 7: Distributions, Testing, and Models}

\sprint{1}{Which Distribution Fits My Data?}{25}
\begin{itemize}[nosep]
  \item Simulate from each: \texttt{np.random.binomial}, \texttt{np.random.poisson}, \texttt{np.random.negative\_binomial}, \texttt{np.random.normal}.
  \item Plot histograms side by side. Then overlay a real dataset (raw RNA-seq counts from Block 5, or a column of allele frequencies or expression values).
  \item Key observation: raw RNA-seq counts are not normal; they are overdispersed and well fit by a negative binomial. That is why DESeq2 and edgeR use negative-binomial models.
\end{itemize}
\begin{winbox}[Checkpoint]
You can look at a histogram of counts and name a plausible distribution family. Save a four-panel comparison figure to your repo.
\end{winbox}

\begin{tipbox}[The distribution cheat sheet]
\begin{itemize}[nosep]
  \item Binomial: counting successes (alt allele reads out of total reads)
  \item Poisson: rare events (mutations per gene)
  \item Negative binomial: counts with extra variation (RNA-seq --- this is the big one)
  \item Normal: continuous measurements after transformation (log-expression, quantitative traits)
  \item Beta: values between 0 and 1 (allele frequencies, methylation)
  \item Hypergeometric: overlap significance (gene set enrichment, Fisher's exact test)
\end{itemize}
\end{tipbox}

\sprint{2}{Multiple Testing: The Genomics Problem}{25}
\begin{itemize}[nosep]
  \item Simulate 1,000,000 null p-values: \texttt{p = np.random.uniform(0, 1, 1\_000\_000)}.
  \item Count how many are less than 0.05. You should see about 50,000: pure noise at the nominal $\alpha$ = 0.05 threshold.
  \item Apply Bonferroni: threshold $=$ $0.05 / 1{,}000{,}000$. Conservative, and misses small real effects.
  \item Apply Benjamini--Hochberg: controls the false discovery rate (expected proportion of false calls \textit{among rejections}), not the family-wise error rate. This is the genomics standard.
  \item Plot a QQ plot of sorted $-\log_{10}(\text{observed } p)$ vs.\ $-\log_{10}(\text{expected } p)$. Deviation above the diagonal is signal (or inflation).
\end{itemize}
\begin{winbox}[Checkpoint]
You simulated a million null p-values, counted false positives at $p < 0.05$, and applied both Bonferroni and BH corrections. Report how many discoveries survive each.
\end{winbox}

\sprint{3}{Linear Models: The GWAS Engine}{25}
\begin{itemize}[nosep]
  \item A GWAS is a linear regression run millions of times: \texttt{phenotype \textasciitilde{} genotype + covariates}, one model per variant.
  \item In Python use \texttt{statsmodels.OLS}; in R use \texttt{lm()}. Fit a model, interpret the genotype coefficient, and check residuals for problems.
  \item Extend with GLMs: logistic regression for case/control (binary outcomes), Poisson or negative-binomial for counts.
  \item Mixed-effects models add random effects for family structure or batch. Conceptual understanding is enough for this sprint; you will use them in the Running Project.
\end{itemize}
\begin{winbox}[Checkpoint]
You fit a single-SNP association test (the core unit of a GWAS) and can read off the effect size and p-value for the genotype term.
\end{winbox}

\begin{warnbox}[Common mistakes to watch for]
\begin{itemize}[nosep]
  \item Using a t-test when you have covariates to control for (use regression instead)
  \item Treating FDR as the probability any specific result is false (it's the expected proportion)
  \item Ignoring population structure in GWAS (this inflates false positives massively)
  \item Applying normal-distribution tests to count data (use NB or nonparametric tests)
\end{itemize}
\end{warnbox}


\section{Block 8: Dimensionality Reduction and Bayesian Ideas}

\sprint{4}{PCA: Finding Hidden Structure}{25}
\begin{itemize}[nosep]
  \item Run PCA on a 1000 Genomes subset (or any genotype matrix) with \texttt{sklearn.decomposition.PCA}.
  \item Plot PC1 vs.\ PC2, colored by population label. Continental ancestry typically emerges as visible clusters.
  \item Make a scree plot (variance explained per PC). The ``elbow'' suggests how many PCs carry signal rather than noise.
  \item This is why GWAS includes the top PCs as covariates: they absorb ancestry and reduce population-structure false positives.
\end{itemize}
\begin{winbox}[Checkpoint]
Your PC1--PC2 plot shows clear clusters by population label, and your scree plot shows a visible elbow (commonly around 3--10 components on 1000 Genomes-scale data).
\end{winbox}

\sprint{5}{UMAP and Clustering for Single-Cell}{25}
\begin{itemize}[nosep]
  \item UMAP is a nonlinear embedding that preserves local structure well and global structure poorly. It is useful for single-cell data, but easy to over-interpret.
  \item Rule to remember: in a UMAP plot, between-cluster distances and orientations are not meaningful; two clusters close together are not necessarily more similar than two far apart.
  \item Leiden clustering is graph-based and has a \texttt{resolution} parameter. Run several resolutions and compare the partitions; ``correct'' resolution is a modeling choice, not a fact.
  \item Hands-on: load a pre-processed single-cell dataset and make UMAP and cluster plots in scanpy or Seurat.
\end{itemize}
\begin{winbox}[Checkpoint]
You produced a UMAP with Leiden clusters overlaid at two different resolutions and wrote one sentence in your journal about what changes and what stays the same.
\end{winbox}

\sprint{6}{Bayesian Thinking (Conceptual)}{25}
\begin{itemize}[nosep]
  \item Bayes' theorem in one line: posterior $\propto$ prior $\times$ likelihood. You update what you believed before seeing the data with what the data tell you.
  \item Empirical Bayes estimates the prior from the data itself. DESeq2's dispersion shrinkage is exactly this: share information across genes to stabilize per-gene estimates.
  \item Why it matters: Bayesian fine-mapping (e.g., SuSiE) reports 95\% credible sets of causal variants instead of individual p-values.
  \item You do not need to implement MCMC. You do need to know what ``95\% credible set'' means: with posterior probability 0.95 under the model, the true causal variant is one of the variants in the set.
\end{itemize}
\begin{winbox}[Checkpoint]
You can explain Bayes' theorem with one genomics example (e.g., DESeq2 dispersion shrinkage or SuSiE credible sets) in plain English, in two to three sentences.
\end{winbox}

\begin{restbox}[End of Blocks 7--8]
This is the most conceptually dense part of the curriculum. Take some time to let it settle before moving into applications.
\end{restbox}

\subsection{Checkpoint:}
\begin{itemize}[nosep]
  \item You can identify which distribution fits a genomic dataset
  \item You can apply and explain Bonferroni and BH correction
  \item You can fit a linear model and interpret coefficients
  \item You can run PCA and explain what population structure looks like
  \item You can explain empirical Bayes at a conceptual level
\end{itemize}

\subsection{Resources}
\begin{itemize}[nosep]
  \item StatQuest YouTube (Josh Starmer) --- highly recommended for visual explanations of statistical concepts
  \item Holmes \& Huber -- \textit{Modern Statistics for Modern Biology} (free online)
  \item Storey \& Tibshirani -- Statistical significance for genomewide studies (the FDR paper)
\end{itemize}
